% GENERATED FILE. Edit canonical lessons, then run npm run generate:books.
% canonical-source-hash: fnv1a64:3e3f28b4ad4ede5c
% canonical-lessons: SA-C130-murkhah, SA-C130-alasah, SA-C130-sabalah, SA-C130-durbalah, SA-C130-svacchah

\chapter{Foolish, Lazy, Strong, Weak, Clear}
\label{ch:sa-foolish-lazy-strong-weak-clear}
\hlchaptermodality{130}

\begin{chapteropening}
\textbf{By the end of this chapter:} \emph{I can say foolish, lazy, strong, weak and clear.}

Complete the last lesson of chapter 130: i can say foolish, lazy, strong, weak and clear.
\end{chapteropening}

\section[mūrkhaḥ]{\sk{मूर्खः} (mūrkhaḥ) — foolish}

\label{lesson:SA-C130-murkhah}



\begin{quote}
Before the new one: say the Sanskrit for brave, then the Sanskrit for clever.
\end{quote}

\subsection*{You'll want to know: \sk{मूर्खः}}
\textbf{\sk{मूर्खः}} — \emph{mūrkhaḥ} — \textquotedblleft{}foolish\textquotedblright{}.

\begin{grammarlens}[title={What you've built}]
One more word for this chapter.
\end{grammarlens}

\subsection*{Your turn}
\begin{itemize}
\raggedright
  \item \emph{Say it:} \emph{mūrkhaḥ}
  \item \emph{Say it:} \emph{mūrkhaḥ}, once more
  \item \emph{Say it:} say \emph{mūrkhaḥ}
  \item \emph{Recall:} say the Sanskrit for brave, then the Sanskrit for clever, then say \emph{mūrkhaḥ} again
\end{itemize}

\begin{tcolorbox}[breakable,colback=teal!4,colframe=teal!35!black,title={Before you move on}]
What does \sk{मूर्खः} mean? (\textquotedblleft{}Foolish\textquotedblright{}.) Say it once more.
\end{tcolorbox}
\section[alasaḥ]{\sk{अलसः} (alasaḥ) — lazy}

\label{lesson:SA-C130-alasah}



\begin{quote}
Before the new one: say the Sanskrit for clever, then the Sanskrit for foolish.
\end{quote}

\subsection*{You'll want to know: \sk{अलसः}}
\textbf{\sk{अलसः}} — \emph{alasaḥ} — \textquotedblleft{}lazy\textquotedblright{}.

\begin{grammarlens}[title={What you've built}]
One more word for this chapter.
\end{grammarlens}

\subsection*{Your turn}
\begin{itemize}
\raggedright
  \item \emph{Say it:} \emph{alasaḥ}
  \item \emph{Say it:} \emph{alasaḥ}, once more
  \item \emph{Say it:} say \emph{alasaḥ}
  \item \emph{Recall:} say the Sanskrit for clever, then the Sanskrit for foolish, then say \emph{alasaḥ} again
\end{itemize}

\begin{tcolorbox}[breakable,colback=teal!4,colframe=teal!35!black,title={Before you move on}]
What does \sk{अलसः} mean? (\textquotedblleft{}Lazy\textquotedblright{}.) Say it once more.
\end{tcolorbox}
\section[sabalaḥ]{\sk{सबलः} (sabalaḥ) — strong}

\label{lesson:SA-C130-sabalah}



\begin{quote}
Before the new one: say the Sanskrit for foolish, then the Sanskrit for lazy.
\end{quote}

\subsection*{You'll want to know: \sk{सबलः}}
\textbf{\sk{सबलः}} — \emph{sabalaḥ} — \textquotedblleft{}strong\textquotedblright{}.

\begin{grammarlens}[title={What you've built}]
One more word for this chapter.
\end{grammarlens}

\subsection*{Your turn}
\begin{itemize}
\raggedright
  \item \emph{Say it:} \emph{sabalaḥ}
  \item \emph{Say it:} \emph{sabalaḥ}, once more
  \item \emph{Say it:} say \emph{sabalaḥ}
  \item \emph{Recall:} say the Sanskrit for foolish, then the Sanskrit for lazy, then say \emph{sabalaḥ} again
\end{itemize}

\begin{tcolorbox}[breakable,colback=teal!4,colframe=teal!35!black,title={Before you move on}]
What does \sk{सबलः} mean? (\textquotedblleft{}Strong\textquotedblright{}.) Say it once more.
\end{tcolorbox}
\section[durbalaḥ]{\sk{दुर्बलः} (durbalaḥ) — weak}

\label{lesson:SA-C130-durbalah}



\begin{quote}
Before the new one: say the Sanskrit for lazy, then the Sanskrit for strong.
\end{quote}

\subsection*{You'll want to know: \sk{दुर्बलः}}
\textbf{\sk{दुर्बलः}} — \emph{durbalaḥ} — \textquotedblleft{}weak\textquotedblright{}.

\begin{grammarlens}[title={What you've built}]
One more word for this chapter.
\end{grammarlens}

\subsection*{Your turn}
\begin{itemize}
\raggedright
  \item \emph{Say it:} \emph{durbalaḥ}
  \item \emph{Say it:} \emph{durbalaḥ}, once more
  \item \emph{Say it:} say \emph{durbalaḥ}
  \item \emph{Recall:} say the Sanskrit for lazy, then the Sanskrit for strong, then say \emph{durbalaḥ} again
\end{itemize}

\begin{tcolorbox}[breakable,colback=teal!4,colframe=teal!35!black,title={Before you move on}]
What does \sk{दुर्बलः} mean? (\textquotedblleft{}Weak\textquotedblright{}.) Say it once more.
\end{tcolorbox}
\section[svacchaḥ]{\sk{स्वच्छः} (svacchaḥ) — clear, clean}

\label{lesson:SA-C130-svacchah}



\begin{quote}
Before the new one: say the Sanskrit for strong, then the Sanskrit for weak.
\end{quote}

\subsection*{You'll want to know: \sk{स्वच्छः}}
\textbf{\sk{स्वच्छः}} — \emph{svacchaḥ} — \textquotedblleft{}clear, clean\textquotedblright{}.

\begin{grammarlens}[title={What you've built}]
One more word for this chapter.
\end{grammarlens}

\subsection*{Your turn}
\begin{itemize}
\raggedright
  \item \emph{Say it:} \emph{svacchaḥ}
  \item \emph{Say it:} \emph{svacchaḥ}, once more
  \item \emph{Say it:} say \emph{svacchaḥ}
  \item \emph{Recall:} say the Sanskrit for strong, then the Sanskrit for weak, then say \emph{svacchaḥ} again
\end{itemize}

\begin{tcolorbox}[breakable,colback=teal!4,colframe=teal!35!black,title={Before you move on}]
What does \sk{स्वच्छः} mean? (\textquotedblleft{}Clear, clean\textquotedblright{}.) Say it once more.
\end{tcolorbox}
